\section{Introduction}
\label{sec:intro}

\aidraft{%
A recurring claim in the applied regime-switching literature is that identifying the market's
latent state and trading on it improves investor outcomes.
Recent work built on statistical jump models \citep{nystrup2020eswa,nystrup2021,shu2024,jmmpc2025}
reports that regime-aware allocations beat buy-and-hold and raise risk-adjusted returns.
The claim is attractive because the estimator is attractive: jump models produce stable,
interpretable state sequences from a handful of return-based features, without the fragility
that made hidden Markov regime models hard to deploy out of sample.%
}
\needcite{Shu--Yu--Mulvey exact headline improvement + their benchmark}
\aidraft{%
The cleanest instance is the allocation sequel of \citet{shu2025dynamic}: regime forecasts enter
expected returns through a gradient-boosted classifier, the covariance is a plain EWMA, and the
benchmarks are minimum-variance, mean-variance, and equal-weight --- with no exposure-matched,
volatility-targeted, or cost-matched control.
\reviewnote{Adam: confirm the 60/20/20-vs-60/40 framing you want here.}%
}

\aidraft{%
We take the estimator seriously as a measurement device and separate two questions the
literature tends to fuse.
First: does the instrument \emph{measure} the market's risk state reliably --- is the label a
stable, reproducible, live-computable quantity?
Second: does \emph{conditioning a decision} on that label add value over a simple estimator that
never names the regime?
Our central finding is that these questions have opposite answers, and that the reported
economic value lives in the gap between them.%
}

\aidraft{%
On the first question the instrument is exceptional.
A $K{=}2$ jump model on return-only downside features produces a label that is \stabJM{} stable
under $\pm 2$-year perturbations of the training window --- where an HMM ensemble on the same
data agrees with itself only \stabHMM{} of the time --- switches a modest \switchesYr{} times per
year, and reproduces to \overlapAgree{} agreement on a \overlapDays{}-day out-of-sample overlap
when carried to the present through a live data splice.
As a market-state sensor it is genuine: it flags \bearsCaught{} ex-post bear markets.
This is a measurement instrument worth having, and Section~\ref{sec:instrument} validates it as such.%
}

\aidraft{%
On the second question every use we tested is dominated by a reactive estimator computing the
decision-relevant quantity directly.
We evaluate four uses under strict preregistered, one-look discipline (Section~\ref{sec:discipline}):
trading the label as an exposure signal (Section~\ref{sec:race1}), conditioning a multi-asset
covariance estimate on it (Section~\ref{sec:race2}), mapping the filter's own evidence margin to
a state probability (Section~\ref{sec:race3}), and rotating a defensive cross-asset sleeve on it
(Section~\ref{sec:race4}).
The exposure strategy's apparent \feeJMBH{}~bps/yr utility fee over buy-and-hold sits inside
every exposure-matched null band and loses a significant \feeJMVT{}~bps/yr to volatility
targeting; covariance conditioning adds a band-interior \feeAmatch{}~bps/yr and loses
\feeBreact{}~bps/yr to a plain EWMA covariance; the evidence margin is beaten by an EWMA-volatility
probability in \dgpCells{} synthetic environments; and cross-asset rotation's drawdown protection
is matched by a volatility target at equal average exposure.
A single mechanism, demonstrated in simulation before any real-data run, unifies them
(Section~\ref{sec:mechanism}): regime value at daily horizons is an information-\emph{timing}
claim, and a causal detector's recognition lag lands it exactly where even a perfectly informed
oracle loses to an estimator that pays no recognition toll.%
}

\aidraft{%
Our contribution is therefore both constructive and deflationary.
We deliver a regime instrument validated to an unusual standard --- perturbation stability, live
operation, and a claim-by-claim ledger of what is out-of-sample-validated versus merely
descriptive (Section~\ref{sec:monitor}, the monitor).
And we show, with exposure-matched nulls and reactive co-primary baselines that the source
literature omits, that its decision value at daily horizons is an artifact of comparison choices.
The instrument-versus-strategy distinction --- usually a footnote --- is the paper's thesis:
measurement validity is not decision value, and honest baselines separate them.%
}
